\chapter{Discussão}
\label{ch:discussao}

Este capítulo interpreta os resultados do Capítulo~\ref{ch:resultados} à luz da fundamentação teórica e organiza os achados nos três artefatos de contribuição anunciados na Seção~\ref{sec:justificativa}: um catálogo de requisitos funcionais e não funcionais (Seção~\ref{sec:disc-requisitos}), um conjunto de problemas recorrentes de usabilidade e acessibilidade (Seção~\ref{sec:disc-ux}) e uma taxonomia das decisões de arquitetura observadas nos casos (Seção~\ref{sec:disc-arquitetura}). Antes disso, a Seção~\ref{sec:disc-adocao} discute o padrão geral de aceitação encontrado entre os frequentadores, e a Seção~\ref{sec:disc-triangulacao} confronta a visão dos gestores com a dos clientes.

%=======================================================================
\section{Aceitação do cardápio digital pelos frequentadores}
\label{sec:disc-adocao}

% Leitura prévia (n = 41, 12/11/2026; refazer com a base final):
% - Preferência: cerca de 78% preferem o cardápio físico e 5% o digital;
%   ainda assim, cerca de 66% acham que o cardápio digital "veio para ficar".
% - EE favorável (56% acham fácil aprender; só 27% relatam dificuldade de
%   leitura/navegação), mas PE baixo (7% acham mais prático que o físico)
%   e BI2 quase nulo (2% recomendariam um lugar por ter cardápio digital).
% - TRU alto (cerca de 71% confiam nos preços), apesar de relatos de itens
%   indisponíveis nas respostas abertas.
% - "Ter os dois" é a resposta aberta mais recorrente na pergunta "Por quê?".
% Argumento sugerido: o cardápio digital é aceito como algo inevitável, mas
% não é desejado; a barreira não é aprender a usar (EE), e sim o valor
% percebido (PE, PV, HM). Dialogar com Venkatesh et al. (2012), que apontam
% PE como preditor mais forte de BI, e com Costa (2022), que também encontrou
% descontentamento em Brasília.

\pendente{discutir o contraste entre facilidade de uso percebida e utilidade percebida, e a preferência por soluções híbridas, relacionando com a UTAUT2 \cite{venkatesh2012consumer}, com os trabalhos relacionados e com a exigência municipal de cardápio físico \cite{camarapoa2024cardapio}}

%=======================================================================
\section{Catálogo de requisitos}
\label{sec:disc-requisitos}

Os requisitos a seguir foram extraídos das respostas abertas do questionário, dos problemas técnicos relatados e das entrevistas com gestores. Cada requisito é classificado como funcional (RF) ou não funcional (RNF), conforme \citeonline{sommerville2016software}, e os não funcionais são associados à característica de qualidade correspondente da ISO/IEC 25010 \cite{iso25010}. A coluna ``Origem'' indica se a necessidade apareceu entre os clientes (C), entre os gestores (G) ou em ambos.

\begin{table}[!ht]
    \caption{Catálogo preliminar de requisitos para soluções de cardápio digital com autoatendimento}
    \label{tab:requisitos}
    \centering%
    \begin{minipage}{.98\textwidth}
        \small
        \renewcommand{\arraystretch}{1.25}
        \begin{tabular*}{\textwidth}{l@{\extracolsep{\fill}}>{\raggedright\arraybackslash}p{7.2cm}>{\raggedright\arraybackslash}p{2.6cm}c}
            \toprule
            \textbf{ID} & \textbf{Requisito} & \textbf{Tipo / ISO 25010} & \textbf{Origem} \\
            \midrule
            R01 & Confirmar ao cliente, de forma visível, que o pedido foi recebido e em que estado ele está & RF & C \\
            R02 & Refletir em tempo real a indisponibilidade de itens & RF & C+G \\
            R03 & Permitir consultar o cardápio e pedir sem cadastro ou login obrigatório & RF & C \\
            R04 & Permitir que várias pessoas da mesma mesa consultem e montem o pedido ao mesmo tempo & RF & C \\
            R05 & Permitir consultar o cardápio antes de chegar ao estabelecimento & RF & C \\
            R06 & Carregar rapidamente em conexões lentas, evitando arquivos pesados como PDF & RNF / eficiência de desempenho & C \\
            R07 & Oferecer alternativa quando não há internet (Wi-Fi do estabelecimento, cardápio físico) & RNF / confiabilidade & C+G \\
            R08 & Navegar entre categorias sem recarregar a página nem perder a posição & RNF / usabilidade & C \\
            R09 & Permitir ajuste de tamanho de fonte e contraste adequado & RNF / usabilidade (acessibilidade) & C \\
            R10 & Exibir muitos itens por tela e permitir comparar opções sem rolagem extensa & RNF / usabilidade & C+G \\
            R11 & Vincular os itens do cardápio ao estoque, por ficha técnica, para atualizar a disponibilidade automaticamente & RF & G \\
            \bottomrule
        \end{tabular*}
        \fonte{Elaborado pelo autor a partir das respostas ao questionário e das entrevistas.}
    \end{minipage}
\end{table}

A entrevista do Caso D, que escolheu não usar cardápio digital, ainda assim contribui para o catálogo. A crítica do gestor à rolagem e à dificuldade de comparar itens coincide com relatos abertos dos frequentadores sobre ``idas e vindas'' entre páginas (R10), e a defesa de manter uma alternativa física para quem não tem internet ou smartphone converge com R07. A atualização imediata de itens esgotados aparece como a principal vantagem reconhecida do digital (R02), e a operação do caso mostra que ela depende da integração com o estoque (R11), funcionalidade que o bar já usava no sistema de gestão, mesmo sem cardápio digital.

\pendente{acrescentar os requisitos dos casos A, B e C e do fornecedor e comentar onde gestores e clientes concordam e onde entram em conflito}

%=======================================================================
\section{Problemas recorrentes de usabilidade e acessibilidade}
\label{sec:disc-ux}

A Tabela~\ref{tab:problemas-ux} relaciona os problemas mais citados pelos frequentadores às heurísticas de \citeonline{nielsen1994heuristics} e, quando aplicável, aos critérios de sucesso das WCAG~2.2 \cite{wcag22}.

% Frequência prévia dos problemas técnicos (n = 41): layout confuso 71%;
% QR Code que não abre 68%; página que não carrega 63%; cadastro ou dados
% pessoais exigidos 49%; tela de tablet/totem que não responde 41%; falta de
% opção de acessibilidade 32%; preço ou informação desatualizados 15%.

\begin{table}[!ht]
    \caption{Problemas recorrentes de usabilidade e acessibilidade em cardápios digitais}
    \label{tab:problemas-ux}
    \centering%
    \begin{minipage}{.98\textwidth}
        \small
        \renewcommand{\arraystretch}{1.25}
        \begin{tabular*}{\textwidth}{>{\raggedright\arraybackslash}p{5.2cm}@{\extracolsep{\fill}}>{\raggedright\arraybackslash}p{4.6cm}>{\raggedright\arraybackslash}p{3.4cm}}
            \toprule
            \textbf{Problema relatado} & \textbf{Heurística de Nielsen} & \textbf{WCAG 2.2} \\
            \midrule
            Layout confuso ou difícil de ler & Estética e design minimalista; consistência e padrões & --- \\
            Cardápio em PDF difícil de ler no celular & Estética e design minimalista & 1.4.10 Reflow \\
            Dúvida se o pedido foi feito ou cobrado & Visibilidade do estado do sistema & --- \\
            Navegação que recarrega ou perde a posição ao voltar & Controle e liberdade do usuário & --- \\
            Necessidade de rolar a tela para comparar itens & Reconhecimento em vez de memorização & --- \\
            Item exibido, mas indisponível & Prevenção de erros; correspondência com o mundo real & --- \\
            Cadastro ou login exigido para ver o cardápio & Flexibilidade e eficiência de uso & --- \\
            Falta de ajuste de fonte e contraste & --- & 1.4.3 Contraste; 1.4.4 Redimensionar texto \\
            \bottomrule
        \end{tabular*}
        \fonte{Elaborado pelo autor a partir das respostas ao questionário, das entrevistas, de \citeonline{nielsen1994heuristics} e de \citeonline{wcag22}.}
    \end{minipage}
\end{table}

\pendente{comentar cada problema com trechos das respostas abertas e das entrevistas; discutir a diferença entre QR Code no celular do cliente e tablet ou totem do estabelecimento, que vários respondentes preferem}

%=======================================================================
\section{Decisões de arquitetura nos casos estudados}
\label{sec:disc-arquitetura}

A taxonomia proposta organiza as decisões de arquitetura em seis dimensões, derivadas da Seção~\ref{sec:arquitetura} e do Bloco~2 do roteiro de entrevista, e posiciona cada caso estudado em cada dimensão \cite{bass2021software,richards2020fundamentals}. A coluna do fornecedor descreve a solução que ele oferece; no perfil D, que não usa cardápio digital, as dimensões são preenchidas com as razões da escolha quando o gestor as mencionou.

\begin{table}[!ht]
    \caption{Taxonomia de decisões de arquitetura observadas nos casos}
    \label{tab:taxonomia-arquitetura}
    \centering%
    \begin{minipage}{.98\textwidth}
        \footnotesize
        \renewcommand{\arraystretch}{1.25}
        \begin{tabular*}{\textwidth}{>{\raggedright\arraybackslash}p{2.8cm}@{\extracolsep{\fill}}>{\raggedright\arraybackslash}p{2.0cm}>{\raggedright\arraybackslash}p{2.0cm}>{\raggedright\arraybackslash}p{2.0cm}>{\raggedright\arraybackslash}p{2.0cm}>{\raggedright\arraybackslash}p{2.0cm}}
            \toprule
            \textbf{Dimensão} & \textbf{Caso A} & \textbf{Caso B} & \textbf{Caso C} & \textbf{Caso D} & \textbf{Fornecedor} \\
            \midrule
            Dispositivo (cliente ou estabelecimento) & & & & Nenhum para o cliente; cardápio impresso & \\
            Sistema próprio ou contratado & & & & Gestão contratada como serviço (Saipos) & \\
            Integração com PDV, estoque e cozinha & & & & Pedido do garçom no sistema, com impressoras, pagamento, estoque e notas fiscais & \\
            Comportamento sem internet & & & & Não abordado & \\
            Formato do cardápio (PDF, web, aplicativo) & & & & Impresso; mesma arte no Instagram & \\
            Identificação do cliente (cadastro, crédito) & & & & Nenhuma & \\
            \bottomrule
        \end{tabular*}
        \fonte{Elaborado pelo autor a partir das entrevistas com gestores e com o fornecedor.}
    \end{minipage}
\end{table}

\pendente{preencher a tabela com os dados das entrevistas e discutir as escolhas entre desenvolver e contratar (build versus buy) em pequenos estabelecimentos}

%=======================================================================
\section{Visão dos gestores e visão dos clientes}
\label{sec:disc-triangulacao}

\pendente{confrontar as categorias das entrevistas com os resultados do questionário, indicando convergências, divergências e complementaridades \cite{venkatesh2013bridging}}
